\section{Discussion}

The central result of this study is that anchor-based transfer improves generalization because it changes the prediction problem itself. Standard drug-target affinity models learn an absolute mapping from a protein--drug pair to a binding score, which is fragile under dataset shift because the model must infer compatibility from representations alone. Our anchor formulation instead asks a relative question: given that anchor protein $a$ is already known to bind drug $d$ strongly, how should query protein $q$ be scored relative to $a$ for the same drug? This is more transferable because it conditions prediction on experimentally grounded binding evidence rather than requiring the model to recover binding from sequence and chemistry alone. The anchor acts as a task-specific reference state, turning prediction from ``does $q$ bind $d$?'' into ``how does $q$ compare with a known binder of $d$?'' This reframing is consistent with findings in the transfer learning literature that relative or relational representations generalize better across domains than absolute features \cite{pan2010transfer,zhuang2021transfer}.

\subsection{What Anchor Transfer Actually Learns}

The anchor-stratified analysis on Davis (Section~\ref{sec:anchor_quartile}) provides direct evidence for what anchor transfer learns. Three observations are particularly revealing.

First, pairwise models (DeepDTA, ESM-DTA) achieve CI~$\approx 0.5$ on Davis---indistinguishable from random---despite strong performance on DTC (CI~$> 0.77$). These models have seen 77\% of Davis proteins during training (with different drugs), yet they cannot transfer binding knowledge to novel compounds even when evaluated with realistic Tanimoto-retrieved anchors. This confirms that pairwise architectures memorize protein--drug co-occurrences rather than learning transferable binding determinants.

Second, the complementary relationship between anchor quality and encoder capacity reveals the mechanism of transfer. V2-650M dominates in Q1--Q2 (weak anchors, pKi 7.0--9.1), achieving CI 0.691 versus 0.583 for V2-35M. But in Q4 (strong anchors, pKi 9.6--10.8), V2-35M overtakes V2-650M (CI 0.744 vs.\ 0.706). This trade-off has a natural interpretation: when the anchor provides a strong binding reference point, even a modest protein encoder can exploit the relative signal---the anchor constrains the prediction space sufficiently that additional encoder capacity yields diminishing returns. When the anchor is weak, the model must rely more heavily on the protein representation to discriminate, and the larger encoder compensates. Anchor quality and encoder capacity are partially substitutable resources for the same underlying task: estimating how a query protein's binding potential compares to that of a known binder.

Third, the per-protein metric distributions (Figure~\ref{fig:davis_per_protein}) confirm that V2-650M's advantage is consistent across proteins, not driven by outliers. The CI distribution is shifted rightward (mean 0.642) compared to baselines that cluster at random (DeepDTA 0.520, ESM-DTA 0.501). This consistency indicates that anchor transfer captures a general property of protein--drug binding rather than exploiting dataset-specific shortcuts.

\subsection{The Role of Protein Structural Similarity}

The success of anchor transfer rests on an implicit assumption: that structurally or functionally similar proteins bind similar drugs with predictable affinity relationships. The model does not receive explicit structural information---it operates on frozen protein language model embeddings---but the anchor mechanism creates an inductive pathway through which structural similarity becomes functionally relevant.

When the model is given an anchor protein that binds a drug strongly, it implicitly learns to compare the query protein's representation against the anchor's. If the two proteins share structural features relevant to binding (conserved active sites, similar fold topology, compatible binding pockets), the model can transfer the anchor's binding evidence to predict the query's affinity. The ESM-2 embeddings capture aspects of protein structure through their pre-training on evolutionary data \cite{lin2023esm2}, and the anchor mechanism leverages this structural signal in a binding-specific way.

This interpretation is supported by the ProstT5 results. V2-ProstT5-BDB achieves the highest global AUROC on Davis (0.768), surpassing V2-650M-DTC (0.757). ProstT5 \cite{heinzinger2023prostt5} is explicitly trained on 3Di structural alphabet representations from Foldseek \cite{van2024fast}, making its embeddings directly structure-aware. The fact that structure-aware embeddings improve anchor transfer performance suggests that the model benefits from better encoding of the structural similarity between anchor and query proteins.

The practical implication is that anchor transfer is most effective when the training set contains proteins that are structurally related to the evaluation targets. DTC-trained models benefit from 77\% protein sequence overlap with Davis, providing the model with anchor proteins from the same kinase family. BDB-trained models, which have broader protein diversity but less kinase enrichment, achieve lower Davis CI (0.602 for V2-35M-BDB vs.\ 0.624 for V2-35M-DTC). However, anchor quality filtering partially compensates: removing anchors with pKi~$< 7$ improves V2-650M-BDB global AUROC from 0.615 to 0.698, demonstrating that anchor quality matters more than anchor quantity.

\subsection{Limitations and the Need for Deeper Models}

Despite the strong results on Davis, several limitations indicate that deeper models and more diverse training data are needed for truly general-purpose drug-target affinity prediction.

\paragraph{Scope of generalization: compound novelty, not protein family novelty.}
We emphasize that our claim is \emph{not} that anchor transfer generalizes to protein families the model has never seen. The DTC training set shares 77\% of its proteins with Davis by sequence, and BDB shares 84\%. What the model has learned is a \emph{general representation of how drugs interact with known protein families}: given a protein that is structurally related to proteins in the training set, and given a known strong binder of a chemically similar drug, the model can predict whether the query protein will bind a novel compound. This is compound novelty---new drugs for known protein families---which is the dominant scenario in practical drug discovery (lead optimization, scaffold hopping, selectivity profiling, repurposing \cite{pushpakom2019drug_repurposing}).

For the model to generalize to entirely novel protein families, it would need training data spanning a broader range of protein structures and binding modes. The contrast between DTC-trained and BDB-trained models illustrates this directly: BDB has greater protein diversity but less kinase enrichment, and its models achieve lower Davis CI (0.602 vs.\ 0.624 for V2-35M). Broader generalization requires more proteins and more drugs---anchor transfer provides the mechanism, but the mechanism's reach is bounded by the structural diversity of the training data.

\paragraph{Drug overlap is deeper than SMILES strings suggest.}
A critical finding of our overlap audit (Table~\ref{tab:overlap_audit}) is that raw SMILES string comparison dramatically underestimates true drug overlap: while 0\% of Davis drugs match DTC by string identity, 83.8\% match by canonical SMILES and 89.7\% by InChIKey. Most Davis kinase inhibitors are literally present in DTC under different string representations. We address this by explicitly excluding all canonical duplicates from the anchor retrieval pool before evaluation, ensuring genuine zero molecular overlap.

After exclusion, the Tanimoto similarity distribution between Davis drugs and their nearest DTC match spans a wide range (12 drugs below 0.4, 22 between 0.4--0.6, 18 between 0.6--0.8). Table~\ref{tab:tanimoto_strat} shows that anchor transfer works even for chemically dissimilar compounds: the [0--0.6) bin (34 drugs with no close DTC analog) achieves CI 0.656 and AUROC 0.651. The best performance occurs at moderate similarity [0.6--0.8) (CI 0.708, AUROC 0.778), where the anchor's drug is structurally related but not identical to the query drug---providing useful binding context without redundancy.

This analysis has broader implications for the DTA literature. Many cross-dataset evaluations report zero drug overlap based on string comparison alone, which we show can mask near-complete molecular overlap. We recommend that future DTA benchmarks verify overlap at the canonical SMILES or InChIKey level to ensure genuine compound novelty.

\paragraph{Shallow architecture limits expressivity.}
The current architecture uses three independent pairwise MLPs (anchor--drug, query--drug, anchor--query) without higher-order interactions. This design is intentionally simple to isolate the contribution of the anchor mechanism from architectural complexity. However, it limits the model's ability to capture complex three-way relationships. Attention-based fusion, graph neural networks over the anchor--query--drug triplet, or deeper cross-attention modules could improve the model's capacity to reason about structural compatibility. The anchor quality--encoder capacity trade-off (Section~\ref{sec:anchor_quartile}) suggests that deeper models would most benefit the weak-anchor regime, where the current architecture already shows the largest gains from scaling ESM-2.

\paragraph{Training data diversity.}
The contrast between DTC-trained and BDB-trained models highlights the importance of training data composition. DTC is kinase-enriched, which helps on the kinase-focused Davis benchmark. BDB is broader but noisier, with 36\% of anchors having pKi~$< 7$ (weak or non-binding). Training on curated, structurally diverse datasets with verified high-quality binding data---such as filtered BindingDB with anchor quality thresholds---is likely to improve generalization to non-kinase families. The improvement from anchor quality filtering (pKi~$\geq 7$) already demonstrates this principle.

\paragraph{Anchor retrieval and future directions.}
The main evaluation uses Tanimoto chemical similarity for anchor retrieval, which requires no privileged access to evaluation data. More sophisticated retrieval strategies---learned retrieval modules that estimate transfer utility, or multi-anchor ensembles that aggregate evidence from several known binders---could further improve performance and provide calibrated uncertainty when anchors disagree. These extensions would be especially valuable for compounds with low Tanimoto similarity to the training set, where the current nearest-neighbor retrieval is weakest.

\subsection{Broader Implications}

These results suggest that cross-dataset DTA should be approached as a \emph{knowledge transfer problem over known interactions}, not solely as a representation learning problem. The failure of ESM-DTA (the strongest same-dataset model, AUROC 0.707 on DTC test) to generalize (CI 0.501 on Davis) demonstrates that better protein representations alone do not solve the transfer problem. The anchor mechanism provides the missing ingredient: a way to condition predictions on experimentally grounded binding evidence, making the model's reasoning explicitly relational rather than implicitly associative.

Critically, the model learns how drug--protein--protein relationships work within known protein families. When the model encounters a query protein that is structurally similar to proteins in its training set, and the query drug has a chemical analog with known binding data, the anchor mechanism enables transfer of that binding knowledge. This is not a limitation but a precise characterization of the method's scope: anchor transfer works because it exploits the structural regularity of protein families to predict binding of novel compounds. Extending this to novel protein families requires broader training data---more proteins from diverse structural classes, paired with more drugs from diverse chemical series. The anchor transfer mechanism itself is family-agnostic; it is the training data that determines which families the model can serve.
